% ==============================================================================
% MASTER'S THESIS — EBERHARD KARLS UNIVERSITÄT TÜBINGEN
% Faculty of Economics and Social Sciences
% Chair of Econometrics (Prof. Dr. Michael Knaus)
%
% Title: Outcome Weights, Kappa Estimators, and Covariate Balance
%        in LATE Estimation
%
% Compile with: pdflatex → biber → pdflatex → pdflatex
%   or in Overleaf: set compiler to pdfLaTeX, bibliography tool to Biber
% ==============================================================================

\documentclass[12pt, a4paper, oneside]{article}

% ------------------------------------------------------------------------------
% ENCODING & LANGUAGE
% ------------------------------------------------------------------------------
\usepackage[utf8]{inputenc}
\usepackage[T1]{fontenc}
\usepackage[english]{babel}
\usepackage{lmodern}
\usepackage{microtype}

% ------------------------------------------------------------------------------
% FONTS  — matches standard econometrics journals (Times-like serif)
% ------------------------------------------------------------------------------

\usepackage[scaled=0.90]{helvet}
\usepackage{courier}

% ------------------------------------------------------------------------------
% PAGE GEOMETRY
% ------------------------------------------------------------------------------
\usepackage[
  a4paper,
  top    = 2.5cm,
  bottom = 2.5cm,
  left   = 3.0cm,
  right  = 2.5cm,
  headheight = 14pt
]{geometry}

% ------------------------------------------------------------------------------
% SPACING
% ------------------------------------------------------------------------------
\usepackage{setspace}
\onehalfspacing                % 1.5 line spacing (standard for master theses)

% ------------------------------------------------------------------------------
% HEADERS & FOOTERS
% ------------------------------------------------------------------------------
\usepackage{fancyhdr}
\pagestyle{fancy}
\fancyhf{}
\fancyhead[L]{\small\textit{Outcome Weights, Kappa Estimators, and Covariate Balance}}
\fancyhead[R]{\small\thepage}
\renewcommand{\headrulewidth}{0.4pt}
\fancypagestyle{plain}{%      % Chapter opening pages
  \fancyhf{}
  \fancyfoot[C]{\small\thepage}
  \renewcommand{\headrulewidth}{0pt}
}

% ------------------------------------------------------------------------------
% MATHEMATICS
% ------------------------------------------------------------------------------
\usepackage{amsmath, amssymb, amsthm}
\usepackage{bm}                % bold math symbols

% Theorem-like environments
\theoremstyle{plain}
\newtheorem{proposition}{Proposition}[section]
\newtheorem{lemma}[proposition]{Lemma}
\newtheorem{corollary}[proposition]{Corollary}
\theoremstyle{definition}
\newtheorem{definition}[proposition]{Definition}
\newtheorem{remark}[proposition]{Remark}

% Convenient math macros
\newcommand{\E}{\mathbb{E}}
\newcommand{\R}{\mathbb{R}}
\newcommand{\N}{\mathbb{N}}
\newcommand{\Prob}{\mathbb{P}}
\newcommand{\indic}{\mathbf{1}}
\newcommand{\tauhat}{\hat{\tau}}
\newcommand{\phat}{\hat{p}}
\newcommand{\omegavec}{\boldsymbol{\omega}}
\newcommand{\kappavec}{\boldsymbol{\kappa}}
\DeclareMathOperator{\Var}{Var}
\DeclareMathOperator{\Cov}{Cov}
\DeclareMathOperator{\plim}{plim}

% ------------------------------------------------------------------------------
% TABLES & FIGURES
% ------------------------------------------------------------------------------
\usepackage{booktabs}          % publication-quality tables
\usepackage{threeparttable}    % notes below tables
\usepackage{tabularx}
\usepackage{multirow}
\usepackage{graphicx}
\usepackage{float}
\usepackage{caption}
\usepackage{subcaption}
\captionsetup{
  font      = small,
  labelfont = bf,
  format    = hang,
  width     = 0.95\textwidth
}

% ------------------------------------------------------------------------------
% HYPERLINKS & CROSS-REFERENCES
% ------------------------------------------------------------------------------
\usepackage[
  colorlinks = true,
  linkcolor  = black,
  citecolor  = black,
  urlcolor   = black
]{hyperref}
\usepackage{cleveref}

% ------------------------------------------------------------------------------
% BIBLIOGRAPHY  — Author-year (AER/econometrics standard)
% ------------------------------------------------------------------------------
\usepackage[round,authoryear]{natbib}


% ------------------------------------------------------------------------------
% MISC
% ------------------------------------------------------------------------------
\usepackage{microtype}         % better typography
\usepackage{parskip}           % space between paragraphs instead of indent
\usepackage{enumitem}
\usepackage{xcolor}
\usepackage{lscape}            % landscape pages for wide tables
\usepackage{rotating}
\usepackage{pdflscape}
\usepackage[normalem]{ulem}    % strikethrough if needed

% Code listings (for appendix R code)
\usepackage{listings}
\lstset{
  basicstyle   = \small\ttfamily,
  breaklines   = true,
  frame        = single,
  language     = R,
  commentstyle = \color{gray},
  keywordstyle = \color{blue!70!black},
  numbers      = left,
  numberstyle  = \tiny\color{gray},
  xleftmargin  = 2em
}

% ------------------------------------------------------------------------------
% PDF METADATA
% ------------------------------------------------------------------------------
\hypersetup{
  pdftitle  = {Outcome Weights, Kappa Estimators, and Covariate Balance in LATE Estimation},
  pdfauthor = {Alexander Unger},
  pdfsubject= {Master's Thesis, University of Tübingen},
}


% ==============================================================================
\begin{document}
% ==============================================================================

% ------------------------------------------------------------------------------
% TITLE PAGE
% ------------------------------------------------------------------------------
\begin{titlepage}
  \thispagestyle{empty}
  \centering

  \vspace*{1.5cm}

  {\large \textsc{Eberhard Karls Universität Tübingen}}\\[0.3cm]
  {\large Faculty of Economics and Social Sciences}\\[0.1cm]
  {\normalsize Department of Economics}\\[0.1cm]
  {\normalsize Chair of Data Science in Economics}

  \vspace{1.5cm}

  \rule{\textwidth}{1.5pt}\\[0.4cm]
  {\LARGE \bfseries Comparing Kappa Weighting and Causal Machine \\[0.3cm]
   Learning Estimators via Weight-Based Diagnostics}\\[0.4cm]
  \rule{\textwidth}{1.5pt}

  \vspace{1.5cm}

  {\large Master's Thesis}\\[0.2cm]
  {\normalsize submitted in partial fulfilment of the requirements for the degree of}\\[0.2cm]
  {\large \textit{Master of Science in Data Science in Business and Economics}}

  \vspace{2.0cm}

  \begin{tabular}{ll}
    \textbf{Author:}     & Alexander Unger \\[0.2cm]
    \textbf{Matriculation number:} & [6946672] \\[0.2cm]
    \textbf{First supervisor:}  & Junior Prof.\ Dr.\ Michael Knaus \\[0.2cm]
    \textbf{Submission date:}   & \today \\
  \end{tabular}

  \vfill

  {\small Tübingen, Germany}

\end{titlepage}

% ------------------------------------------------------------------------------
% FRONT MATTER
% ------------------------------------------------------------------------------
\pagenumbering{roman}

% Abstract
\newpage
\section*{Abstract}
\addcontentsline{toc}{section}{Abstract}

This thesis applies the outcome-weights framework of \citet{knaus2024} to kappa-based
local average treatment effect (LATE) estimators of \citet{abadie2003} and
\citet{suw2025}. Any estimator that can be written as a weighted sum of observed outcomes,
$\hat\tau = \sum_i \omega_i Y_i$, admits direct covariate balance diagnostics via Love
plots and effective sample size (ESS) calculations. Knaus (2024) derives such outcome
weights for double machine learning (DML) and generalised random forest estimators;
Appendix A.4 sketches but does not implement the same for kappa estimators. We fill
this gap by deriving closed-form outcome weights for all five kappa-based LATE estimators
and showing that the sum-to-zero condition $\sum_i \omega_i = 0$ is the exact finite-sample
criterion for translation invariance---the property that the LATE estimate is unaffected
by additive rescaling of the outcome variable.

We apply the unified outcome-weight diagnostics to three empirical settings: the Vietnam
draft lottery \citep{angrist1990}, college proximity as an instrument for education
\citep{card1995}, and sibling sex composition as an instrument for fertility
\citep{angrist1998}. A striking finding across applications is that in low-dimensional,
near-random-assignment designs, ESS and the percentage of negative weights are virtually
identical across all estimators and all machine learning learners---a pattern we term
\emph{design dominates learner}. Normalization, not covariate adjustment, drives estimate
divergence. Unnormalized kappa estimators fail translation invariance dramatically in the
Angrist--Evans design, with estimates changing sign under outcome recoding. Normalized
estimators $\hat\tau_u$ and $\hat\tau_{a,10}$ are robust throughout.

\bigskip
\noindent\textbf{Keywords:} local average treatment effect, instrumental variables,
kappa weighting, outcome weights, translation invariance, covariate balance,
double machine learning.

\bigskip
\noindent\textbf{JEL classification:} C14, C21, C26.

% Table of contents
\newpage
\tableofcontents

% List of figures
\newpage
\listoffigures
\addcontentsline{toc}{section}{List of Figures}

% List of tables
\newpage
\listoftables
\addcontentsline{toc}{section}{List of Tables}

% ------------------------------------------------------------------------------
% MAIN TEXT
% ------------------------------------------------------------------------------
\newpage
\pagenumbering{arabic}
\setcounter{page}{1}

% ==============================================================================
\section{Introduction}
\label{sec:intro}
% ==============================================================================

\subsection{Motivation}
\label{subsec:motivation}

% [WRITE HERE — hook: cents vs. dollars failure, ~0.5 page]

\subsection{Background}
\label{subsec:background}

% [WRITE HERE — 2SLS doesn't recover LATE, kappa as alternative, ~0.5 page]

\subsection{Research Gap and Contribution}
\label{subsec:contribution}

% [WRITE HERE — Knaus Appendix A.4 gap, list 4–5 contributions, ~1 page]

\subsection{Road Map}
\label{subsec:roadmap}

% [WRITE HERE — 1 paragraph, chapter-by-chapter]


% ==============================================================================
\section{Econometric Framework}
\label{sec:framework}
% ==============================================================================

\subsection{IV, LATE, and Compliers}
\label{subsec:late}

% [WRITE HERE — potential outcomes, four types, LATE definition, ~1.5 pages]

\subsection{Abadie's Kappa Theorem}
\label{subsec:kappa_theorem}

% [WRITE HERE — Lemma 2.1, three kappa weights, Table 1 of SUW, ~1.5 pages]

\subsection{Kappa-Based LATE Estimators}
\label{subsec:estimators}

% [WRITE HERE — five estimators, MLE vs CBPS, ~1.5 pages]

\subsection{Why Normalisation Matters}
\label{subsec:normalisation}

% [WRITE HERE — TI definition, Prop 3.2 SUW, cents/dollars example, ~1.5 pages]

\subsection{Double Machine Learning and Wald-AIPW}
\label{subsec:dml}

% [WRITE HERE — DML framework, cross-fitting, PLR-IV and Wald-AIPW, ~1.5 pages]

\subsection{The PIVE Framework and Outcome Weights}
\label{subsec:pive}

% [WRITE HERE — Knaus 2024 PIVE structure, normalization classes, ~1 page]

\subsection{Covariate Balance Diagnostics}
\label{subsec:balance}

% [WRITE HERE — SMD, Love plots, ESS, sum-to-zero as TI criterion, ~1 page]


% ==============================================================================
\newpage
% ==============================================================================
% \section{Connecting the Frameworks}
\label{sec:connection}


\section{Connecting the Frameworks}
\label{sec:connection}
% ==============================================================================
 
The preceding chapter introduced two literatures that are closely related but not automatically equivalent. On the one hand, Abadie's kappa theorem and the weighting estimators studied by \citet{sloczynski2025abadie} use functions of $(Z_i,D_i,p(X_i))$ to identify moments of the complier distribution. On the other hand, the PIVE framework of \citet{knaus2024treatment} characterizes estimators through outcome weights $\omega_i$ satisfying $
\hat\tau = \sum_{i=1}^N \omega_i Y_i ,
$
which makes it possible to study covariate balance, effective sample size, negative weights, and normalization properties directly from the finite-sample estimator.

This chapter connects these two perspectives. The connection is not immediate because Abadie's kappa weights and Knaus's outcome weights are different objects. Kappa weights are used to establish identification: they transform population expectations into complier moments. Outcome weights are estimator-specific finite-sample objects: they describe how a particular estimate is assembled from the observed outcomes. Therefore, the fact that a LATE estimator is based on kappa weights does not by itself imply that the kappa weights are the relevant weights for balance diagnostics.

The purpose of this chapter is to derive the relevant outcome weights for the kappa estimators considered in this thesis and to study their normalization properties. Appendix~A.4 of \citet{knaus2024treatment} places several kappa estimators in the PIVE framework. The estimator $\hat\tau_u$ of \citet{uysal2011three}, however, requires separate treatment because it applies a Hájek-type normalization to the instrument-specific inverse-probability weights before forming the Wald ratio. Since $\hat\tau_u$ is the recommended normalized estimator in \citet{sloczynski2025abadie} and is central to the empirical comparisons below, its outcome-weight representation is essential for interpreting the diagnostics used in the applications.

The main result is that both $\hat\tau_u$ and $\hat\tau_{a,10}$ admit exact finite-sample outcome-weight representations. Moreover, their weights satisfy the normalization conditions that imply translation invariance. This shows that the robustness of the normalized kappa estimators to additive outcome transformations is not merely an asymptotic property, but follows directly from their finite-sample weight structure.

\subsection{Kappa Weights versus Outcome Weights}
\label{subsec:distinction}

Two types of weights appear throughout this thesis. They are based on similar ingredients, but they answer different questions.\\
Kappa weights are identification weights. The functions $\kappa_i$, $\kappa_{i1}$, and $\kappa_{i0}$ introduced by \citet{abadie2003semiparametric} are functions of the observed instrument, treatment, and instrument propensity score. Their role is to express moments of the complier distribution in terms of observable population moments. In particular, for suitable functions $g$,
$
\mathbb{E}!\left[\kappa_i g(Y_i,D_i)\right]
$
recovers the corresponding complier moment up to the complier share. Thus, kappa weighting is an identification device: it tells us how the observed population must be reweighted to isolate the complier subpopulation.

Outcome weights have a different role. Following \citet{knaus2024treatment}, outcome weights are sample-level scalars $\omega_i$ such that a given estimator can be written exactly as
$
\hat\tau = \sum_{i=1}^N \omega_i Y_i .
$
These weights are not universal features of the research design. They depend on the algebraic form of the estimator. Consequently, $\hat\tau_u$, $\hat\tau_{a,10}$, and $\hat\tau_a$ generally induce different outcome weights, even though all three are constructed from the same underlying kappa functions.

This distinction matters because the diagnostics used in the empirical part of the thesis are diagnostics of outcome weights, not of kappa weights. Love plots, standardized mean differences, effective sample size, negative weight shares, and the sum-to-zero condition for translation invariance all refer to $\omega_i$. The task is therefore not simply to compute $\kappa_i$, but to derive the estimator-specific mapping from kappa quantities to outcome weights.



% ==============================================================================
\subsection{The PIVE Framework}
\label{subsec:pive_recap}
 
The derivations in this chapter use the outcome-weight formula for
pseudo-IV estimators established by \citet{knaus2024treatment}.
A pseudo-IV estimator (PIVE) is any estimator that solves the empirical moment condition
\begin{equation}
  \mathbb{E}_N\!\left[\widetilde{Z}_i\left(\widetilde{Y}_i
  - \hat\tau\,\widetilde{D}_i\right)\right] = 0,
  \label{eq:pive_moment}
\end{equation}
for scalar pseudo-quantities $\widetilde{Y}_i$, $\widetilde{D}_i$, $\widetilde{Z}_i$
that may depend on nuisance estimates but not on $\hat\tau$ itself.
In matrix notation, \eqref{eq:pive_moment} has the unique solution
\begin{equation}
  \hat\tau
  = \frac{\widetilde{\mathbf{Z}}'\widetilde{\mathbf{Y}}}
         {\widetilde{\mathbf{Z}}'\widetilde{\mathbf{D}}}.
  \label{eq:pive_sol}
\end{equation}
 
The key result is that whenever the pseudo-outcome can be written as a linear
transformation of the original outcome vector, $\widetilde{\mathbf{Y}} = \mathbf{T}\mathbf{Y}$,
the estimator admits an exact outcome-weight representation.
 
\begin{proposition}[Outcome weights of a PIVE, {\citealt[Proposition~1]{knaus2024treatment}}]
\label{prop:pive_weights}
Suppose $\widetilde{\mathbf{Z}}'\widetilde{\mathbf{D}} \neq 0$ and there exists an
$N \times N$ matrix $\mathbf{T}$ such that $\widetilde{\mathbf{Y}} = \mathbf{T}\mathbf{Y}$.
Then $\hat\tau = \boldsymbol{\omega}'\mathbf{Y}$ with
\begin{equation}
  \boldsymbol{\omega}'
  = \left(\widetilde{\mathbf{Z}}'\widetilde{\mathbf{D}}\right)^{-1}
    \widetilde{\mathbf{Z}}'\mathbf{T}.
  \label{eq:pive_weights}
\end{equation}
\end{proposition}
 
For kappa estimators, the pseudo-outcome is $\widetilde{Y}_i = t_i Y_i$ for some
observation-specific scalar $t_i$ that depends only on $(Z_i, D_i, \hat{p}_i)$.
The transformation matrix $\mathbf{T}$ is therefore diagonal, and its existence is
immediate: no outcome regression or smoother matrix is involved.
This means the conditions in \citet[Section~4.2]{knaus2024treatment} governing
smoother-based estimators such as DML or causal forests do not apply here.
The only operative condition is the separate normalization of the instrument-group
weights, which is specific to $\hat\tau_u$ and is addressed directly in
Section~\ref{subsec:omega_u}.
 
 
% ==============================================================================
\subsection{Outcome Weights for the Uysal Estimator}
\label{subsec:omega_u}
% ==============================================================================

This section derives the outcome weights of the normalized estimator $\hat\tau_u$.
The estimator can be written as a Wald ratio in which both the reduced form and the
first stage are computed using normalized inverse-probability weights for the two
instrument groups. This is the kappa-estimator analogue of the normalized Wald-IPW
case discussed by \citet{knaus2024treatment}.

First, define the two instrument-group normalizing constants as
\begin{equation}
  s_1 = \sum_{j=1}^{N} \frac{Z_j}{\hat{p}_j},
  \qquad
  s_0 = \sum_{j=1}^{N} \frac{1-Z_j}{1-\hat{p}_j}.
  \label{eq:u_s1_s0}
\end{equation}

The corresponding normalized instrument inverse-probability weights are
\begin{equation}
  \lambda_{1i}^{z,n} = \frac{N Z_i/\hat{p}_i}{s_1},
  \qquad
  \lambda_{0i}^{z,n} = \frac{N(1-Z_i)/(1-\hat{p}_i)}{s_0}.
  \label{eq:u_norm_lambda}
\end{equation}

By construction, both sets of weights sum to $N$:
\begin{equation}
  \sum_{i=1}^{N} \lambda_{1i}^{z,n} = N,
  \qquad
  \sum_{i=1}^{N} \lambda_{0i}^{z,n} = N.
  \label{eq:u_norm_lambda_sums}
\end{equation}

Now define the diagonal transformation matrix
\begin{equation}
  \mathbf{T}^u = \operatorname{diag}\!\left(
    \lambda_{11}^{z,n} - \lambda_{01}^{z,n},\;
    \ldots,\;
    \lambda_{1N}^{z,n} - \lambda_{0N}^{z,n}
  \right),
  \label{eq:u_T_matrix}
\end{equation}
so that the $i$th diagonal element is
\begin{equation}
  T^u_{ii} = \lambda_{1i}^{z,n} - \lambda_{0i}^{z,n}.
  \label{eq:u_T_diag}
\end{equation}

With this notation, the Uysal estimator can be written as
\begin{equation}
  \hat\tau_u
  = \frac{N^{-1}\mathbf{1}_N'\mathbf{T}^u\mathbf{Y}}
         {N^{-1}\mathbf{1}_N'\mathbf{T}^u\mathbf{D}}.
  \label{eq:u_wald_matrix}
\end{equation}

This is a PIVE representation with
\begin{equation}
  \widetilde{\mathbf{Y}} = \mathbf{T}^u\mathbf{Y},
  \qquad
  \widetilde{\mathbf{D}} = \mathbf{T}^u\mathbf{D},
  \qquad
  \widetilde{\mathbf{Z}} = \mathbf{1}_N.
  \label{eq:u_pive_components}
\end{equation}

The denominator of the Wald ratio is the normalized IPW first stage. Define it as
\begin{equation}
  \widehat\Delta_D = N^{-1}\mathbf{1}_N'\mathbf{T}^u\mathbf{D},
  \label{eq:u_delta_matrix}
\end{equation}
which in scalar form reads
\begin{equation}
  \widehat\Delta_D
  = \frac{\sum_{j=1}^{N} Z_j D_j/\hat{p}_j}{s_1}
  - \frac{\sum_{j=1}^{N} (1-Z_j)D_j/(1-\hat{p}_j)}{s_0}.
  \label{eq:u_delta_scalar}
\end{equation}

Applying the outcome-weight formula for PIVEs (Proposition~\ref{prop:pive_weights})
yields
\begin{equation}
  \boldsymbol{\omega}^{u\prime}
  = \left(N^{-1}\mathbf{1}_N'\mathbf{T}^u\mathbf{D}\right)^{-1}
    N^{-1}\mathbf{1}_N'\mathbf{T}^u
  = \widehat\Delta_D^{-1}\,N^{-1}\mathbf{1}_N'\mathbf{T}^u.
  \label{eq:u_omega_matrix}
\end{equation}

Therefore, the outcome weight of observation $i$ is
\begin{equation}
  \omega_i^u
  = \frac{1}{N\widehat\Delta_D}
    \left(\lambda_{1i}^{z,n} - \lambda_{0i}^{z,n}\right).
  \label{eq:u_omega_lambda}
\end{equation}

Substituting the normalized weights from \eqref{eq:u_norm_lambda} gives the
closed-form expression
\begin{equation}
  \boxed{
  \omega_i^u
  = \frac{1}{\widehat\Delta_D}
    \left[
      \frac{Z_i}{\hat{p}_i\, s_1}
      - \frac{1-Z_i}{(1-\hat{p}_i)\,s_0}
    \right].
  }
  \label{eq:u_omega_scalar}
\end{equation}

Each outcome weight is the difference between observation $i$'s normalized $Z{=}1$
and $Z{=}0$ inverse-probability weights, rescaled by the estimated first stage
$\widehat\Delta_D$.

% ------------------------------------------------------------------
\paragraph{Normalization.}
% ------------------------------------------------------------------

We verify the three conditions for full normalization in \citet{knaus2024treatment}.

\smallskip
\noindent\textit{(i) Weights sum to zero.}
Using \eqref{eq:u_omega_matrix},
\begin{equation}
  \sum_{i=1}^{N} \omega_i^u
  = \boldsymbol{\omega}^{u\prime}\mathbf{1}_N
  = \widehat\Delta_D^{-1}\,N^{-1}\,\mathbf{1}_N'\mathbf{T}^u\mathbf{1}_N.
  \label{eq:u_sum_start}
\end{equation}
Since $\mathbf{T}^u$ is diagonal with entries
$\lambda_{1i}^{z,n} - \lambda_{0i}^{z,n}$,
\begin{equation}
  \mathbf{1}_N'\mathbf{T}^u\mathbf{1}_N
  = \sum_{i=1}^{N}\left(\lambda_{1i}^{z,n} - \lambda_{0i}^{z,n}\right)
  = N - N = 0,
  \label{eq:u_T_ones_zero}
\end{equation}
where the last equality uses \eqref{eq:u_norm_lambda_sums}. Hence,
\begin{equation}
  \sum_{i=1}^{N} \omega_i^u = 0. \qquad\checkmark
  \label{eq:u_sum_zero}
\end{equation}

\smallskip
\noindent\textit{(ii) Treated weights sum to one.}
\begin{equation}
  \sum_{i=1}^{N} \omega_i^u D_i
  = \boldsymbol{\omega}^{u\prime}\mathbf{D}
  = \widehat\Delta_D^{-1}\,N^{-1}\,\mathbf{1}_N'\mathbf{T}^u\mathbf{D}
  = \widehat\Delta_D^{-1}\,\widehat\Delta_D
  = 1. \qquad\checkmark
  \label{eq:u_treated_sum_one}
\end{equation}

\smallskip
\noindent\textit{(iii) Untreated weights sum to minus one.}
This follows directly from (i) and (ii):
\begin{equation}
  \sum_{i=1}^{N} \omega_i^u(1-D_i)
  = \sum_{i=1}^{N} \omega_i^u - \sum_{i=1}^{N} \omega_i^u D_i
  = 0 - 1
  = -1. \qquad\checkmark
  \label{eq:u_untreated_sum}
\end{equation}

\medskip
Collecting the three conditions:
\begin{equation}
  \sum_{i=1}^{N} \omega_i^u = 0,
  \qquad
  \sum_{i=1}^{N} \omega_i^u D_i = 1,
  \qquad
  \sum_{i=1}^{N} \omega_i^u(1-D_i) = -1.
  \label{eq:u_full_normalization}
\end{equation}

Therefore, $\hat\tau_u$ is fully normalized in the sense of \citet{knaus2024treatment}.
The key step is the separate normalization of the instrument-group inverse-probability
weights in \eqref{eq:u_norm_lambda}, which forces both groups to carry equal total
mass $N$. Their difference therefore sums to zero \eqref{eq:u_T_ones_zero}, which
propagates directly to $\sum_i\omega_i^u = 0$. The treated-weight normalization
follows from the Wald denominator itself: $\widehat\Delta_D$ is defined as exactly the
quantity $N^{-1}\mathbf{1}_N'\mathbf{T}^u\mathbf{D}$ that appears in the numerator of
\eqref{eq:u_treated_sum_one}, so the cancellation is exact in every finite sample.
\begin{remark}[The role of Hájek normalization: contrast with $\hat\tau_{a,1}$]
\label{rem:hajek_contrast}
The sum-to-zero condition $\sum_i \omega_i^u = 0$ rests entirely on
equation~\eqref{eq:u_T_ones_zero}, which holds because both instrument groups carry
equal total mass $N$ after normalization. This makes transparent why the structurally
similar estimator $\hat\tau_{a,1}$ of \citet{frolich2007nonparametric} and \citet{tan2006regression} fails
translation invariance. That estimator uses the same Wald-IPW form as $\hat\tau_u$
but without the within-group normalization: the raw weights $\lambda_{1i} = Z_i/\hat{p}_i$
and $\lambda_{0i} = (1-Z_i)/(1-\hat{p}_i)$ enter directly, so the two group totals are
\begin{equation}
  \sum_{i=1}^N \lambda_{1i} = s_1,
  \qquad
  \sum_{i=1}^N \lambda_{0i} = s_0,
  \label{eq:unnorm_sums}
\end{equation}
which differ whenever the instrument does not assign equal probability across the sample.
The analogue of equation~\eqref{eq:u_T_ones_zero} then becomes
\begin{equation}
  \sum_{i=1}^N \left(\lambda_{1i} - \lambda_{0i}\right) = s_1 - s_0 \neq 0,
  \label{eq:unnorm_diff}
\end{equation}
and the resulting weight sum satisfies
\begin{equation}
  \sum_{i=1}^N \omega_i^{a,1}
  = \widehat\Delta_D^{a,1\,-1} N^{-1}(s_1 - s_0) \neq 0.
  \label{eq:unnorm_weight_sum}
\end{equation}
By the translation-invariance identity $\hat\tau(Y + k\mathbf{1}_N) - \hat\tau(Y)
= k\sum_i\omega_i$, a non-zero weight sum means that adding a constant $k$ to every
outcome shifts the estimate by $k \cdot \widehat\Delta_D^{a,1\,-1} N^{-1}(s_1 - s_0)$.
The Hájek normalization that defines $\hat\tau_u$ relative to $\hat\tau_{a,1}$ is
therefore not a cosmetic refinement but the single algebraic step that determines
whether the estimator is translation invariant.
This is the theoretical counterpart of the empirical finding in
Section~\ref{sec:vietnam}, where $\hat\tau_{a,1}$ and $\hat\tau_{a,0}$ produce
estimates that change sign under outcome recoding while $\hat\tau_u$ does not.
\end{remark}
 

% ---------------------------------------------------------------
\newpage
 
% ==============================================================================

% ==============================================================================
\subsection{Normalization Classification and Summary}
\label{subsec:classification}
% ==============================================================================

\subsubsection{Two Notions of Normalization}
\label{sssec:two_norms}

The preceding sections use the word normalization in two senses that are related but
logically distinct, and it is worth being explicit about both before presenting the
classification.\\
In the kappa-estimator literature, an estimator is called \emph{normalized} if the
inverse-probability or kappa weights used in its construction are rescaled so that the
denominator of the estimator converges to the true complier share
\citep{sloczynski2025abadie}. For $\hat\tau_u$, this means the instrument-group weights
for $Z=1$ and $Z=0$ are separately divided by their group totals before forming the
Wald ratio. For $\hat\tau_{a,10}$, it means the $\kappa_1$- and $\kappa_0$-weighted
sums are each divided by their own total. The unnormalized estimators $\hat\tau_a$,
$\hat\tau_{a,1}$, and $\hat\tau_{a,0}$ omit one or both of these rescaling steps.
This is an internal property of how the estimator is constructed.\\
In the outcome-weight framework of \citet{knaus2024treatment}, normalization refers to
a property of the final weights $\omega_i$ placed on the observed outcomes. An estimator
is \emph{fully normalized} in this sense if
\begin{equation}
  \sum_{i=1}^N \omega_i = 0,
  \qquad
  \sum_{i=1}^N \omega_i D_i = 1,
  \qquad
  \sum_{i=1}^N \omega_i (1-D_i) = -1.
  \label{eq:full_norm_conditions}
\end{equation}
These conditions say that the outcome weights behave like a normalized
treated-minus-control contrast: the treated side carries total mass $+1$ and the
untreated side carries total mass $-1$, so the overall sum is zero.

The first condition in \eqref{eq:full_norm_conditions} has a direct causal
interpretation through a simple algebraic identity. Suppose the outcome variable is
shifted by a constant $k \in \mathbb{R}$, replacing $Y_i$ with $Y_i + k$ for every
observation. This corresponds, for instance, to expressing log wages in cents rather
than dollars, or coding a binary labour-force-participation variable as $\{1,2\}$
rather than $\{0,1\}$. The constant $k$ carries no information about the causal
effect of the treatment. For any estimator $\hat\tau = \sum_i \omega_i Y_i$, the
shifted estimate equals
\begin{equation}
  \hat\tau(Y + k)
  = \sum_{i=1}^N \omega_i (Y_i + k)
  = \sum_{i=1}^N \omega_i Y_i + k \sum_{i=1}^N \omega_i
  = \hat\tau(Y) + k \sum_{i=1}^N \omega_i.
  \label{eq:ti_algebra}
\end{equation}
Rearranging gives the translation-invariance identity:
\begin{equation}
  \hat\tau(Y + k) - \hat\tau(Y) = k \sum_{i=1}^N \omega_i.
  \label{eq:ti_identity}
\end{equation}
This identity holds exactly in finite samples, for any realization of the data and
any propensity score estimate. The right-hand side vanishes if and only if
$\sum_i \omega_i = 0$. 
Thus, the sum-to-zero condition is the necessary and sufficient finite-sample
condition for the estimator to be invariant to additive outcome transformations.
Full normalization in the Knaus sense is stronger: it additionally requires the
treated and untreated components of the outcome weights to sum to $+1$ and $-1$,
respectively.
When $\sum_i \omega_i \neq 0$, the
estimate changes by $k \sum_i \omega_i$ for every unit shift in $k$: an estimator
applied to log wages in cents will produce a numerically different LATE than the
same estimator applied to log wages in dollars, even though the two outcomes contain
identical information about the causal effect. \citet{sloczynski2025abadie} call this
property translation invariance and show in Proposition~3.2 that it holds for
$\hat\tau_u$ and $\hat\tau_{a,10}$ but fails for the unnormalized estimators. The
derivations in this chapter show that the same property is an exact finite-sample
consequence of the outcome-weight structure derived in
Sections~\ref{sssec:derive_a10} and~\ref{sssec:derive_u}.

\subsubsection{How the Two Notions Are Connected}
\label{sssec:norm_connection}

For the estimators studied in this thesis, the two normalization concepts co-occur:
every estimator that is normalized in the SUW sense turns out to be fully normalized
in the Knaus sense, and every unnormalized estimator fails both conditions. This is not
a coincidence. As shown in Sections~\ref{sssec:derive_a10} and~\ref{sssec:derive_u},
the same algebraic operation that makes $\hat\tau_u$ and $\hat\tau_{a,10}$ normalized
in the SUW sense is also the step that forces $\sum_i \omega_i = 0$.\\
For $\hat\tau_{a,10}$, separate rescaling of the $\kappa_1$ and $\kappa_0$ weighted sums
forces each to sum to one, so their difference sums to zero unconditionally. For
$\hat\tau_u$, the within-group Hájek normalization forces $\sum_i \lambda_{1i}^{z,n} =
\sum_i \lambda_{0i}^{z,n} = N$, so the difference of the two group totals is $N - N = 0$,
which propagates directly to $\sum_i \omega_i^u = 0$.\\
On the other hand the unnormalized estimators fail both conditions for the same reason in reverse.
$\hat\tau_a$ uses the raw kappa contrast $\kappa_{i1} - \kappa_{i0}$ without separate
rescaling, giving $\sum_i \omega_i^a = (\sum_i \kappa_{i1} - \sum_i \kappa_{i0}) /
\sum_i \kappa_i$, which is not zero in general \citep[Appendix~A.4]{knaus2024treatment}.
$\hat\tau_{a,1}$ uses the raw IPW contrast $\lambda_{1i} - \lambda_{0i}$ without
Hájek correction, giving $\sum_i \omega_i^{a,1} = \widehat\Delta_D^{a,1\,-1} N^{-1}
(s_1 - s_0) \neq 0$ whenever $s_1 \neq s_0$.\\
It is important to note that this co-occurrence is specific to the kappa family. The two
normalization concepts are logically independent: it is possible in principle to construct
an estimator that is normalized in the SUW sense but not in the Knaus sense, or vice
versa. The alignment observed here reflects the particular algebraic structure of
kappa-based Wald estimators, not a general theorem.

\subsubsection{Classification Table}
\label{sssec:class_table}

Table~\ref{tab:normalization} classifies all estimators studied in this thesis according
to both normalization criteria, following the protocol of \citet[Table~4--5]{knaus2024treatment}.

\begin{table}[t]
\centering
\caption{Normalization classification of kappa and DML estimators}
\label{tab:normalization}
\begin{threeparttable}
\renewcommand{\arraystretch}{1.3}
\begin{tabular}{@{}lllcccc@{}}
\toprule
Estimator & & SUW norm. & $\sum_i\omega_i = 0$ & $\sum_i\omega_iD_i$ & $\sum_i\omega_i(1-D_i)$ & Knaus class \\
\midrule
$\hat\tau_u^{cb}$ & (CBPS) & yes & \checkmark\,(exact) & $= +1$    & $= -1$    & Fully norm. \\
$\hat\tau_u^{ml}$ & (MLE)  & yes & \checkmark\,(exact) & $= +1$    & $= -1$    & Fully norm. \\
$\hat\tau_{a,10}$ &        & yes & \checkmark\,(exact) & $= +1$    & $= -1$    & Fully norm. \\[4pt]
$\hat\tau_a$      &        & no  & $\times$            & $\neq +1$ & $\neq -1$ & Fully unnorm. \\
$\hat\tau_{a,1}$ & & no & $\times$ & $= +1$ & $\neq -1$ & Untreated-unnorm. \\
$\hat\tau_{a,0}$  &        & no  & $\times$            & $\neq +1$ & $= -1$    & Treated-unnorm. \\[4pt]
Wald-AIPW & (DML) & n/a & \checkmark\,(approx.) & $\approx +1$ & $\approx -1$ & Scale-norm. (std.)$^{a}$ \\
\bottomrule
\end{tabular}
\begin{tablenotes}
  \small
  \item \textit{SUW norm.} indicates whether the estimator is normalized in the sense
  of \citet{sloczynski2025abadie}, i.e.\ whether the kappa or IPW weights are
  separately rescaled before forming the Wald ratio.
  \textit{Knaus class} follows \citet[Table~4--5]{knaus2024treatment}.
  \checkmark\,(exact): the condition holds for all data realizations by algebraic
  construction; \checkmark\,(approx.): holds approximately in finite samples subject
  to Conditions~3 and~5 of \citet{knaus2024treatment}, which conflict with standard
  flexible implementations. $\hat\tau_{a,0}$: treated-unnormalized because
  $\sum_i\omega_i^{a,0}(1-D_i) = -1$ holds by kappa orthogonality but
  $\sum_i\omega_i^{a,0} \neq 0$ \citep[Appendix~A.4]{knaus2024treatment}.
  $^{a}$ Wald-AIPW is scale-normalized in standard implementations because
$\sum_i\omega_i = 0$ under Condition~3 of \citet{knaus2024treatment}.
Full normalization additionally requires Condition~5b, which holds only if the same
smoother predicts both outcome and treatment — conflicting with standard practice
of using separate models for each nuisance parameter
\citep[Section~4.2.3]{knaus2024treatment}.
\end{tablenotes}
\end{threeparttable}
\end{table}

Three observations from Table~\ref{tab:normalization} deserve emphasis.\\
First, $\hat\tau_u$ and $\hat\tau_{a,10}$ achieve Knaus full normalization by different
algebraic mechanisms, even though both are SUW-normalized. $\hat\tau_{a,10}$ does so by
separately rescaling the $\kappa_1$ and $\kappa_0$ weighted sums; the sum-to-zero
follows unconditionally from the structure of the kappa weights alone. $\hat\tau_u$ does
so via Hájek normalization of the instrument-group IPW weights; the sum-to-zero follows
from the equal-mass property \eqref{eq:u_norm_lambda_sums}.\\
Second, Wald-AIPW achieves Knaus full normalization only approximately in finite samples,
and only under Conditions~3 and~5 of \citet{knaus2024treatment}, which conflict with
standard flexible implementations such as gradient-boosted trees. The kappa estimators
$\hat\tau_u$ and $\hat\tau_{a,10}$ achieve it exactly and unconditionally, which is a
stronger guarantee.\\
Third, the practical consequence of the $\sum_i\omega_i = 0$ column is given directly
by identity~\eqref{eq:ti_identity}: estimators with $\sum_i\omega_i \neq 0$ will produce
different LATE estimates when the outcome variable is shifted by a constant. The
empirical diagnostics in Chapters~\ref{sec:vietnam}--\ref{sec:card_ae} compute
$\sum_i\omega_i$ numerically for every estimator and every dataset, and verify the
sum-to-zero condition against machine precision. The function
\texttt{kappa\_outcome\_weights()} in \texttt{functions\_kappa.R} implements the
closed-form expressions derived in Sections~\ref{sssec:derive_a10}
and~\ref{sssec:derive_u} and the identity $\hat\tau = \sum_i\omega_i Y_i$ is verified
via \texttt{check\_weight\_identity()} for every specification reported.





% ==============================================================================
\section{Empirical Application: Military Service and Wages}
\label{sec:vietnam}
% ==============================================================================

\subsection{Data and Design}
\label{subsec:vietnam_data}

% [WRITE HERE — SIPP 1984, N=3027, Z=draft lottery, one-sided noncompliance, ~1 page]

\subsection{Point Estimates and Replication}
\label{subsec:vietnam_estimates}

% [TABLE: SUW Table 2 replication — normalised vs unnormalised, cents vs dollars]

% Placeholder for Table 2
\begin{table}[H]
  \centering
  \caption{Kappa estimators: log wages (Vietnam draft lottery)}
  \label{tab:vietnam_estimates}
  \begin{threeparttable}
    \begin{tabular}{llccc}
      \toprule
      & & \multicolumn{3}{c}{Covariate specification} \\
      \cmidrule(lr){3-5}
      Panel & Estimator & Linear age & Cubic age & Saturated age \\
      \midrule
      \multicolumn{5}{l}{\textit{A. 2SLS benchmark}} \\
      & 2SLS & & & \\[4pt]
      \multicolumn{5}{l}{\textit{B. Normalised estimates}} \\
      & $\hat\tau_u^{cb}$ & & & \\
      & $\hat\tau_u^{ml}$ & & & \\
      & $\hat\tau_{a,10}$ & & & \\[4pt]
      \multicolumn{5}{l}{\textit{C. Unnormalised estimates}} \\
      & $\hat\tau_a$       & & & \\
      & $\hat\tau_t$       & & & \\
      & $\hat\tau_{a,0}$   & & & \\
      \bottomrule
    \end{tabular}
    \begin{tablenotes}
      \small\item Robust standard errors in parentheses.
      $^{*}p<0.10$, $^{**}p<0.05$, $^{***}p<0.01$.
      [ADD NOTES ON DATA AND REPLICATION OF SUW TABLE 2]
    \end{tablenotes}
  \end{threeparttable}
\end{table}

\subsection{Double Machine Learning Comparison}
\label{subsec:vietnam_dml}

% [WRITE HERE — PLR-IV, Wald-AIPW, three learners, design dominates learner, ~2 pages]

\subsection{Outcome Weight Diagnostics and Covariate Balance}
\label{subsec:vietnam_weights}

% [WRITE HERE — weight diag table, Love plots, ESS=5, 54% negative weights, ~2 pages]

% Placeholder for Table 3
\begin{table}[H]
  \centering
  \caption{Outcome weight diagnostics: Vietnam (cubic specification)}
  \label{tab:vietnam_weights}
  \begin{threeparttable}
    \begin{tabular}{lrrrr}
      \toprule
      Estimator & $\sum_i\omega_i$ & ESS & \% neg. & Max $|\omega_i|$ \\
      \midrule
      PLR-IV       (DML)             & & & & \\
      Wald-AIPW    (DML)             & & & & \\
      $\hat\tau_u^{cb}$ (kappa, CBPS) & & & & \\
      $\hat\tau_u^{ml}$ (kappa, MLE)  & & & & \\
      $\hat\tau_{a,10}$ (kappa)        & & & & \\
      $\hat\tau_a$      (kappa)        & & & & \\
      $\hat\tau_t$      (kappa)        & & & & \\
      $\hat\tau_{a,0}$  (kappa)        & & & & \\
      \bottomrule
    \end{tabular}
    \begin{tablenotes}
      \small\item $\sum_i\omega_i = 0$ iff translation invariant (Proposition~\ref{}).
      ESS $= 1/\sum_i\omega_i^2$.
    \end{tablenotes}
  \end{threeparttable}
\end{table}

% Placeholder for Figure 1 (Love plots)
\begin{figure}[H]
  \centering
  % \includegraphics[width=\textwidth]{figures/vietnam_love_plots.pdf}
  \caption{Love plots: outcome-weight-adjusted covariate balance, Vietnam (cubic specification)}
  \label{fig:vietnam_love}
  \medskip
  \small [REPLACE WITH ACTUAL FIGURE]
\end{figure}


% ==============================================================================
\section{Empirical Applications: Card (1995) and Angrist--Evans (1998)}
\label{sec:card_ae}
% ==============================================================================

\subsection{Card (1995): College Proximity and Returns to Education}
\label{subsec:card}

\subsubsection{Data and Design}

% [WRITE HERE — NLSY, N, Z=proximity, two treatments, two covariate specs, ~0.5 page]

\subsubsection{Point Estimates and Translation Invariance}

% [TABLE: Card replication, four spec × treatment columns]

\subsubsection{Outcome Weight Diagnostics}

% [WRITE HERE — ESS diverges across specs, Kitagawa ESS ≈ 1, ~1 page]

\subsection{Angrist and Evans (1998): Childbearing and Labour Supply}
\label{subsec:ae}

\subsubsection{Data and Design}

% [WRITE HERE — Census, N≈254k, Z=samesex, D=morekids, LFP + log income, ~0.5 page]

\subsubsection{Point Estimates and Translation Invariance}

% [TABLE: AE replication, LFP and income codings]

\subsubsection{Outcome Weight Diagnostics and DML Limitation}

% [WRITE HERE — kappa diagnostics, DML infeasible (memory ceiling), ~0.5 page]


% ==============================================================================
\section{Discussion}
\label{sec:discussion}
% ==============================================================================

\subsection{Cross-Application Summary}
\label{subsec:cross_application}

% Placeholder for Table 5 (cross-application comparison)
\begin{table}[H]
  \centering
  \caption{Cross-application comparison of outcome weight diagnostics}
  \label{tab:cross_app}
  \begin{threeparttable}
    \begin{tabular}{llrrrr}
      \toprule
      Application & Estimator & $\sum_i\omega_i$ & ESS & \% neg. & Balance \\
      \midrule
      \multicolumn{6}{l}{\textit{Angrist (1990) — Vietnam (cubic spec.)}} \\
      & DML Wald-AIPW & & & & \\
      & $\hat\tau_u^{cb}$ & & & & \\[4pt]
      \multicolumn{6}{l}{\textit{Card (1995) — Card spec, somecol}} \\
      & DML Wald-AIPW & & & & \\
      & $\hat\tau_u^{cb}$ & & & & \\[4pt]
      \multicolumn{6}{l}{\textit{Card (1995) — Kitagawa spec, somecol}} \\
      & DML Wald-AIPW & & & & \\
      & $\hat\tau_u^{cb}$ & & & & \\[4pt]
      \multicolumn{6}{l}{\textit{Angrist \& Evans (1998) — LFP}} \\
      & $\hat\tau_u^{cb}$ & & & & \\
      \bottomrule
    \end{tabular}
    \begin{tablenotes}
      \small\item Balance = maximum absolute SMD from Love plot.
      DML not reported for Angrist--Evans due to memory constraints
      ($N \approx 254{,}000$; see Section~\ref{subsec:ae}).
    \end{tablenotes}
  \end{threeparttable}
\end{table}

\subsection{The Outcome-Weights Lens: What It Adds}
\label{subsec:lens}

% [WRITE HERE — what Love plots reveal beyond point estimates, ~0.5 page]

\subsection{DML Learner Comparison and Implications}
\label{subsec:dml_learners}

% [WRITE HERE — learner convergence, XGBoost issue, ~0.5 page]

\subsection{Practical Guidance}
\label{subsec:practical}

% [WRITE HERE — when to use which estimator, always check sum(omega), ~0.5 page]

\subsection{Limitations}
\label{subsec:limitations}

% [WRITE HERE — DML memory ceiling for large N, Love plots descriptive not inferential, ~0.3 page]


% ==============================================================================
\section{Conclusion}
\label{sec:conclusion}
% ==============================================================================

% [WRITE HERE — summary of findings, recommendation (tau_u^cb), extensions, ~1.5 pages]


% ==============================================================================
% REFERENCES
% ==============================================================================
\newpage


% ==============================================================================
% APPENDICES
% ==============================================================================
\newpage
\appendix
\renewcommand{\thesection}{\Alph{section}}

\section{Proofs}
\label{app:proofs}

% [Formal proofs of sum-to-zero for tau_u and tau_a10 if needed]

\section{Additional Tables and Figures}
\label{app:additional}

% [Additional Love plots, full weight diagnostics tables, learner comparison details]

\section{Data Sources and Variable Definitions}
\label{app:data}

% [Detailed variable definitions for all three applications]

\section{R Code Overview}
\label{app:code}

Key functions implementing the kappa outcome weights are contained in
\texttt{functions\_kappa.R}. The core function \texttt{kappa\_outcome\_weights()} computes
closed-form $\omega_i$ vectors for all five estimators:

\begin{lstlisting}[caption={Outcome weight construction (functions\_kappa.R, \S k)},
                   label={lst:kappa_ow}]
kappa_outcome_weights <- function(Z, D, p) {
  n  <- length(Z)
  kw <- kappa_weights(Z, D, p)

  # tau_u (Uysal 2011)
  s1  <- sum(Z / p);  s0 <- sum((1 - Z) / (1 - p))
  dD  <- sum(D * Z / p) / s1 - sum(D * (1 - Z) / (1 - p)) / s0
  w_u <- (Z / p / s1 - (1 - Z) / (1 - p) / s0) / dD

  # tau_a10 (Abadie-Cattaneo)
  w_a10 <- kw$kappa1 / sum(kw$kappa1) - kw$kappa0 / sum(kw$kappa0)

  # Unnormalised (common numerator)
  num_w <- (Z - p) / (p * (1 - p)) / n

  list(w_u   = as.vector(w_u),
       w_a10 = as.vector(w_a10),
       w_a   = as.vector(num_w / mean(kw$kappa)),
       w_a1  = as.vector(num_w / mean(kw$kappa1)),
       w_a0  = as.vector(num_w / mean(kw$kappa0)))
}
\end{lstlisting}

All replication code is available at: \texttt{[YOUR GITHUB/CLOUD LINK]}


% ==============================================================================
% SELBSTÄNDIGKEITSERKLÄRUNG (statutory declaration — required by Tübingen)
% ==============================================================================


% ==============================================================================
\bibliographystyle{apalike}
\bibliography{references}

\end{document}
% ==============================================================================
